\documentclass[11pt,a4paper]{article}

\usepackage{amsmath,amssymb,amsthm}
\usepackage[utf8]{inputenc}
\usepackage{hyperref}
\usepackage[margin=1in]{geometry}
\usepackage{xcolor}
\usepackage{booktabs}
\usepackage{array}
\usepackage{float}
\usepackage{mathtools}

\newtheorem{theorem}{Theorem}[section]
\newtheorem{conjecture}[theorem]{Conjecture}
\newtheorem{corollary}[theorem]{Corollary}
\newtheorem{lemma}[theorem]{Lemma}
\newtheorem{proposition}[theorem]{Proposition}
\newtheorem{definition}[theorem]{Definition}
\theoremstyle{remark}
\newtheorem{remark}[theorem]{Remark}
\newtheorem*{remark*}{Remark}

\newcommand{\F}{\mathbb{F}}
\newcommand{\Z}{\mathbb{Z}}
\newcommand{\Sp}{\mathrm{Sp}}
\newcommand{\rad}{\mathrm{rad}}
\newcommand{\rank}{\mathrm{rank}}
\newcommand{\Nanti}{N_{\mathrm{anti}}}

\title{$\mathrm{H}^3$ Opens at $n\ge5$: The Cohomological Origin\\
       of the Mermin Modulus\\
       {\large (Paper XX)}}

\author{Zhou-Li Chen \\ co-nlang Research}
\date{June 2026}

\begin{document}
\maketitle

\begin{abstract}
Paper~XIX established that the fiber of the anticommutation count $\Nanti$
over fixed symplectic rank is a \emph{modulus}: no relative-position invariant
of arity~$\le4$ classifies it.  This paper explains why.

We equip a proper $K_5$ configuration of Lagrangians with a simplicial cochain
complex on the 4-simplex boundary $\partial\Delta^4\cong S^3$.  The Maslov bit
$\mu(a,b,c)\in\F_2$ (whether $\omega\not\equiv0$ on the
triple-intersection data) assembles into a 2-cochain, and the per-matching
anticommutation counts assemble into a 3-cochain $\mathbf{n}$.
Our main results are:
\begin{enumerate}
\item \textbf{(The $n=4$ cocycle theorem.)}  At $n=4$,
  $\mathbf{n}=\delta\mu$ in $C^3(\partial\Delta^4;\F_2)$: the obstruction is
  \emph{exact}, $[\mathbf{n}]=0$ in $H^3(S^3;\F_2)$, and $\Nanti=10$ is even.
  This is the cohomological form of the Paper~XVIII rigidity.
\item \textbf{(The unique vanishing dimension.)}  $[\mathbf{n}]=0$ universally
  \emph{if and only if} $n=4$.  At $n=3$, $\mu\equiv0$ forces $\delta\mu=0$,
  yet all-Fano proper $K_5$ have $\Nanti=15$ so $[\mathbf{n}]=1\ne0$.
  At $n\ge5$, $[\mathbf{n}]$ is generically nontrivial.
  Thus $\Nanti\bmod2=\langle[\mathbf{n}],[S^3]\rangle$ is a genuine H$^3$
  invariant outside the Paper~XVIII dimension.
\item \textbf{(Rank-parity.)}  The symmetric form $B$ on the projected
  Lagrangian satisfies $\rank B=n-3$, giving: $\mu\equiv1$ when $n$ is even;
  $\mu\equiv0$ at $n=3$; $\mu$ varies at $n\ge5$ odd.
\end{enumerate}
The failure of arity-$\le4$ classifiers in Paper~XIX is a direct consequence:
arity-3 data (the $\mu$ values) determine only $\delta\mu$, which equals
$\mathbf{n}$ at $n=4$ but not in general; the natural arity-4 invariant $q_4$
lives in a different cochain complex and is saturated ($99.5\%$).
The modulus \emph{is} the H$^3$ class.
\end{abstract}

\section{Introduction}

A proper $K_5$ is the symplectic skeleton of a Mermin-type contextuality
proof: five Lagrangian subspaces of $\F_2^{2n}$, pairwise meeting in a single
ray, with the ten rays carrying the anticommutation data a Kochen--Specker
argument exploits.  A single integer, the cross-context anticommutation count
$\Nanti$, measures how much of that data is ``twisted.''  The arc of Papers
XVIII--XX is the story of one number and the dimension at which it stops being
forced.

Paper~XVIII~\cite{PaperXVIII} proved that in $\Sp(8,\F_2)$ ($n=4$ qubits)
every proper $K_5$ satisfies $\Nanti=10$ universally---a rigidity forced by
two independent symplectic dimension coincidences (the B0, B1, B2 lemmas).  Paper~XIX~\cite{PaperXIX} showed that
at $n\ge5$ this rigidity softens into a \emph{modulus}: within a fixed
symplectic-rank fiber, $\Nanti\bmod2$ is not determined by any
relative-position invariant of arity~$\le4$.  The modulus witness---two
configurations sharing all Maslov-triple and arity-4 data yet with opposite
$\Nanti$ parities---proves the classification problem is genuinely hard.

Paper~XIX established \emph{that} the modulus exists; it did not explain
\emph{why} low-arity invariants must fail.  That is the purpose of this paper.

\subsection*{The answer in one sentence}
$\Nanti\bmod2$ is the evaluation of the 3-cochain $\mathbf{n}$ against the
fundamental class of $\partial\Delta^4\cong S^3$, and $[\mathbf{n}]=0$
universally if and only if $n=4$: the cohomological obstruction is exact
precisely at the Paper~XVIII dimension, not in general.

\subsection*{Structure of the paper}
Section~\ref{sec:setup} sets up the simplicial cochain complex.
Section~\ref{sec:n4theorem} proves the $n=4$ cocycle theorem.
Section~\ref{sec:rankparity} proves the rank-parity lemma and the periodicity
corollary.
Section~\ref{sec:h3opens} establishes that H$^3$ opens at $n\ge5$.
Section~\ref{sec:obstruction} states the obstruction theorem and explains the
failure of arity-$\le4$ classifiers.
Section~\ref{sec:connections} places the result in the series arc.

\paragraph{Conventions.}
$V=\F_2^{2n}$ with standard symplectic form $\omega$.  A \emph{proper $K_5$}
is a 5-tuple of Lagrangians $(L_0,\ldots,L_4)$ with $\dim(L_i\cap L_j)=1$
for all $i\ne j$; we write $v_{ij}=L_i\cap L_j$.

\section{The simplicial cochain complex}\label{sec:setup}

Everything about a proper $K_5$ is indexed by subsets of the five labels:
pairs of Lagrangians give rays, triples give Maslov data, and the matchings
that carry $\Nanti$ are indexed by the omitted fifth label.  This is precisely
the face structure of a 4-simplex, so the invariants organise into cochains on
it.  The point of this section is to make that organisation explicit; the two
sections after it then compute, at $n=4$, that the degree-3 cochain carrying
$\Nanti$ is the coboundary of the degree-2 Maslov cochain.

\subsection{The 4-simplex boundary}

Label the five Lagrangians by $I=\{0,1,2,3,4\}$.  The standard 4-simplex
$\Delta^4$ has vertex set $I$; its boundary $\partial\Delta^4\cong S^3$ carries
the $\F_2$-cochain complex
\[
  0\to C^0 \xrightarrow{\,\delta\,} C^1 \xrightarrow{\,\delta\,} C^2
    \xrightarrow{\,\delta\,} C^3 \to 0
\]
where $C^k$ is the free $\F_2$-module on $k$-faces:
$|C^0|=5$, $|C^1|=10$, $|C^2|=10$, $|C^3|=5$ (tetrahedra).
The cohomology is $H^0=H^3=\F_2$, $H^1=H^2=0$, and the
fundamental class $[S^3]\in H_3(\partial\Delta^4;\F_2)$ is represented by the
sum of all five tetrahedra.

\subsection{The Maslov 2-cochain}

For a triangular face $\{a,b,c\}\subset I$ define
\begin{equation}\label{eq:maslov}
  \mu(a,b,c) \;=\; \bigl[\,\omega\not\equiv 0 \text{ on }
    D_{abc}\,\bigr] \;\in\;\F_2,
  \qquad
  D_{abc}=\{x\in L_a,\,y\in L_b : x+y\in L_c\}.
\end{equation}
Concretely, $\mu(a,b,c)$ records whether the three Lagrangians sit in
\emph{generic} or in \emph{special} relative position: it is the single bit
that distinguishes a ``twisted'' triple from an untwisted one, and it is the
mod-2 shadow of the classical Maslov triple index.
The assignment $\{a,b,c\}\mapsto\mu(a,b,c)$ defines a 2-cochain
$\mu\in C^2(\partial\Delta^4;\F_2)$.
As shown in Paper~XIX~\cite{PaperXIX}, $\mu$ is $S_3$-symmetric and
$\Sp$-invariant.

\subsection{The anticommutation 3-cochain}

The 15 cross-context pairs $\{v_{ij},v_{k\ell}\}$ (with
$\{i,j\}\cap\{k,\ell\}=\emptyset$) decompose uniquely by their excluded index:
\[
  M_m \;=\; \bigl\{\{v_{ij},v_{k\ell}\} :
    \{i,j,k,\ell\}=I\setminus\{m\}\bigr\}
\]
is the set of the three cross-context pairs whose four indices exhaust
$I\setminus\{m\}$.  The sets $M_0,\ldots,M_4$ partition all 15 pairs.
Define
\[
  n_m \;=\; \#\{\text{anticommuting pairs in }M_m\} \bmod 2 \;\in\;\F_2,
\]
so $\mathbf{n}=(n_0,\ldots,n_4)\in C^3(\partial\Delta^4;\F_2)$.
Since the $M_m$ partition all cross-context pairs,
\begin{equation}\label{eq:nanti}
  \Nanti \bmod 2 \;=\; \sum_{m\in I} n_m
  \;=\; \langle\mathbf{n},\,[S^3]\rangle.
\end{equation}

\subsection{The coboundary of $\mu$}

The coboundary $\delta\mu\in C^3$ evaluates on tetrahedron $m$ as
\[
  (\delta\mu)_m \;=\;
  \sum_{\{a,b,c\}\subset I\setminus\{m\}} \mu(a,b,c)
  \quad(\text{sum over the four triangular faces of tetrahedron }m).
\]
Note that $\sum_m(\delta\mu)_m=\sum_{\{a,b,c\}}\mu(a,b,c)\cdot 2\equiv0$
(each triangle is a face of exactly two tetrahedra), consistent with
$H^3(\partial\Delta^4;\F_2)=\F_2$.

\begin{remark*}
The relevant complex is the \emph{Maslov--Wall relative-position complex}
on the abstract index simplex, not the geometric \v{C}ech intersection nerve.
For a proper $K_5$, triple intersections $L_i\cap L_j\cap L_k=0$, so the
intersection nerve is the 1-skeleton $K_5$ (no 2-cells), which has
$H^3=0$ and cannot carry the class.
\end{remark*}

\section{The $n=4$ cocycle theorem}\label{sec:n4theorem}

At the Paper~XVIII dimension the two cochains collapse for entirely independent
reasons---and they collapse to the same thing, each identically zero.  The
anticommutation cochain $\mathbf{n}$ vanishes because every matching carries an
even count (a transversal-geometry fact, B1+B2), while $\delta\mu$ vanishes
because $\mu$ is the constant cochain $1$ (a rank-parity fact, proved in the
next section).  Their coincidence is what makes $\Nanti=10$ rigid; expressed
cohomologically, it says the obstruction class is exact.

\begin{theorem}[$\mathbf{n}=\delta\mu$ at $n=4$]\label{thm:n4cocycle}
For every proper $K_5$ in $\Sp(8,\F_2)$, $\mathbf{n}=\delta\mu$ in
$C^3(\partial\Delta^4;\F_2)$.  In particular $[\mathbf{n}]=0$ in $H^3(S^3;\F_2)$
and $\Nanti\bmod2=0$.
\end{theorem}

\begin{proof}
We show $n_m\equiv0$ and $(\delta\mu)_m\equiv0$ for every $m$.

\medskip
\noindent\textbf{Part (A): $n_m\equiv0$.}
Each $M_m$ is a perfect matching on the four-vertex set $I\setminus\{m\}$.
By Paper~XVIII's B1 and B2 lemmas~\cite{PaperXVIII}, every proper $K_4$
sub-configuration at $n=4$ has exactly one commuting cross-context pair per
matching, so exactly two of the three pairs in $M_m$ anticommute.
Hence $n_m=2\equiv0\pmod2$.

\medskip
\noindent\textbf{Part (B): $(\delta\mu)_m\equiv0$.}
Theorem~\ref{thm:rankparity} below proves $\mu\equiv1$ at $n=4$.
Therefore each of the four triangular faces of tetrahedron $m$ contributes $1$,
giving $(\delta\mu)_m=1+1+1+1=0\pmod2$.
\end{proof}

\begin{corollary}
At $n=4$ the obstruction is \emph{exact}: $\delta\mu=0$ (trivial coboundary),
$[\mathbf{n}]=0$ in $H^3(S^3;\F_2)=\F_2$, and $\Nanti=10$ is even.
The universality of Paper~XVIII is the statement that the H$^3$ class vanishes.
\end{corollary}

The two parts have different characters: Part~(A) is a consequence of the
transversal geometry (B1+B2) and holds at the level of each matching
independently; Part~(B) is the rank-parity phenomenon proved next.

\section{Rank parity and periodicity}\label{sec:rankparity}

\begin{theorem}[Rank parity]\label{thm:rankparity}
Let $(L_a,L_b,L_c)$ be a proper triple of Lagrangians in $\Sp(2n,\F_2)$.
The symmetric form $B$ on the projected Lagrangian (defined in the proof)
satisfies $\rank B=n-3$, with the following consequences:
\begin{enumerate}
\item[(i)] $n$ \emph{even}: $\mu(a,b,c)=1$ for every proper triple.
\item[(ii)] $n=3$: $\mu(a,b,c)=0$ for every proper triple.
\item[(iii)] $n\ge5$ \emph{odd}: $\mu(a,b,c)$ varies with the triple.
\end{enumerate}
\end{theorem}

\begin{proof}
\emph{Strategy.}  We collapse the pair $(L_a,L_b)$ to a \emph{transverse} pair
of Lagrangians in a symplectic space two dimensions smaller.  There the Maslov
bit $\mu$ becomes the diagonal of an explicit symmetric form $B$ on the third
Lagrangian, and the whole question reduces to a single number: the rank of $B$.
We compute its radical exactly---it is spanned by the two relevant rays---so
$\rank B=n-3$, and the parity of that rank decides all three cases.

\medskip
\noindent\textbf{Reduction to a smaller symplectic space.}
Set $W=L_a+L_b$ and $W^\perp=L_a\cap L_b=\langle v_{ab}\rangle$.
Since $\dim W=2n-1$ and $W^\perp\subset W$ has dimension~1, $W$ is coisotropic.
The quotient $\bar W=W/\langle v_{ab}\rangle$ is symplectic of dimension $2(n-1)$,
with induced form $\bar\omega$.

The images $\bar L_a$ and $\bar L_b$ are transverse Lagrangians of $\bar W$
(dimension $n-1$ each, $\bar L_a\cap\bar L_b=0$), giving a perfect pairing
$\bar\omega:\bar L_a\times\bar L_b\to\F_2$.

\medskip
\noindent\textbf{The projected Lagrangian $\bar L_c$.}
Since $v_{ab}\notin L_c$, the map $L_c\to\bar W$ is injective on
$L_c\cap W$.  If $L_c\subset W$ then $W^\perp=\langle v_{ab}\rangle\subset
L_c^\perp=L_c$ (as $L_c$ is Lagrangian), giving $v_{ab}\in L_a\cap L_c=
\langle v_{ac}\rangle$, contradicting $v_{ab}\ne v_{ac}$ in a proper~$K_5$.
Hence $L_c\not\subset W$, forcing $\dim(L_c+W)=2n$ and
$\dim(L_c\cap W)=n-1$.
Set $\bar L_c=(L_c\cap W)/\langle v_{ab}\rangle$; this is a Lagrangian of
$\bar W$ with $\bar L_c\cap\bar L_a=\langle\bar v_{ac}\rangle$ and
$\bar L_c\cap\bar L_b=\langle\bar v_{bc}\rangle$ (each dimension~1).

\medskip
\noindent\textbf{The symmetric form.}
Decompose $\bar W=\bar L_a\oplus\bar L_b$ and let $p_a,p_b$ be the associated
projections.  Define $B:\bar L_c\times\bar L_c\to\F_2$ by
$B(\bar z,\bar z')=\bar\omega(p_a\bar z,\,p_b\bar z')$.
Since $\bar L_c$ is isotropic, $B$ is symmetric.
The diagonal $q(\bar z)=B(\bar z,\bar z)$ satisfies
$\mu(a,b,c)=[q\not\equiv0]$ (by definition~\eqref{eq:maslov} after reducing
to $\bar W$).
Over $\F_2$, $q\equiv0$ if and only if $B$ is alternating.

\medskip
\noindent\textbf{The radical of $B$.}
We show $\rad B=\langle\bar v_{ac},\bar v_{bc}\rangle$.
A vector $\bar z\in\bar L_c$ lies in $\rad B$ iff
$\bar\omega(p_a\bar z,\,p_b\bar z')=0$ for all $\bar z'\in\bar L_c$,
i.e., $p_a\bar z$ lies in the $\bar\omega$-annihilator (in $\bar L_a$) of
$p_b(\bar L_c)$.
Since $\ker(p_b)=\bar L_a$ meets $\bar L_c$ in
$\bar L_c\cap\bar L_a=\langle\bar v_{ac}\rangle$ (dimension~1), the image
$p_b(\bar L_c)$ has dimension $n-2$, and its annihilator in $\bar L_a$
under the perfect pairing has dimension~1.
As $\bar v_{ac}\in\bar L_c\cap\bar L_a$ is isotropic,
$\bar\omega(\bar v_{ac},p_b\bar z')=0$ for all $\bar z'\in\bar L_c$
(since $\bar v_{ac}\in\bar L_a$ and $\bar L_a$ isotropic), so
the annihilator equals $\langle\bar v_{ac}\rangle$.
By symmetry $p_b\bar z\in\langle\bar v_{bc}\rangle$ as well.
Hence $\rad B=\langle\bar v_{ac},\bar v_{bc}\rangle$.

\medskip
\noindent\textbf{Rank and parity.}
Therefore $\rank B=(n-1)-2=n-3$.  Over $\F_2$, $q\equiv0\iff B$ alternating;
alternating forms have even rank.
\begin{itemize}
\item $n$ \emph{even}: $\rank B=n-3$ is odd, so $B$ is not alternating,
  $q\not\equiv0$, $\mu=1$.\quad(i)
\item $n=3$: $\rank B=0$, so $B$ is the zero form (trivially alternating),
  $q\equiv0$, $\mu=0$.\quad(ii)
\item $n\ge5$ \emph{odd}: $\rank B=n-3\ge2$ is even; the non-degenerate part
  of $B$ on $\bar L_c/\rad B$ may or may not be alternating, so $\mu$
  varies.\quad(iii)
\end{itemize}
Case~(iii) is the whole source of the $n\ge5$ phenomenon.  Once $\rank B$ is
even and positive, $B$ can fall on either side of the alternating/non-alternating
divide, so $\mu$ is no longer pinned by the proper-triple combinatorics---and
with it neither is $\delta\mu$, nor the class $[\mathbf{n}]$ it controls.
\qed
\end{proof}

\begin{corollary}[Rank-parity summary]\label{cor:periodicity}
$\mu\equiv1$ when $n$ is even; $\mu\equiv0$ at $n=3$; $\mu$ varies when
$n\ge5$ is odd.
\begin{center}
\begin{tabular}{ccl}
\toprule
$n$ & $\mu$ uniform? & Observation \\
\midrule
$3$ (odd)  & Yes ($\equiv0$) & $\rank B=0$; by Thm~\ref{thm:rankparity}(ii) \\
$4$ (even) & Yes ($\equiv1$) & $300/300$ proper triples \\
$5$ (odd)  & No (varies)    & $230/300$ triples have $\mu=1$ \\
$6$ (even) & Yes ($\equiv1$) & $200/200$ proper triples \\
\bottomrule
\end{tabular}
\end{center}
\vspace{-2pt}
\noindent
Computational rows verified by \texttt{mu\_rank\_parity.py}~\cite{PaperXXsuppl}.
\end{corollary}

\section{H$^3$ opens at $n\ge5$}\label{sec:h3opens}

Once $\mu$ stops being constant, the coincidence that pinned $\mathbf{n}=\delta\mu$
at $n=4$ breaks: $\mu$ ceases to be a cocycle, $\delta\mu\ne0$, and $\mathbf{n}$
need no longer agree with it.  The residue $\mathbf{n}-\delta\mu$ is exactly the
nontrivial $H^3$ class, and it shows up as an odd $\Nanti$.

\begin{proposition}\label{prop:h3opens}
For $n=5$, generically $\mathbf{n}\ne\delta\mu$, and the class
$[\mathbf{n}]\in H^3(S^3;\F_2)$ is nontrivial.
\end{proposition}

\begin{proof}[Proof (numerical)]
By Corollary~\ref{cor:periodicity}, $\mu$ varies at $n=5$.  Over $480$ proper
$K_5$ configurations, some $(\delta\mu)_m\ne0$ in $419$; and over a disjoint
run of $360$ configurations (seeds $1000+7s$), $\mathbf{n}\ne\delta\mu$ in
$336$.
Scripts: \texttt{nerve\_cochain.py}, \texttt{n4\_cocycle.py}~\cite{PaperXXsuppl}.
\end{proof}

The opening of H$^3$ is sharp:

\begin{center}
\begin{tabular}{lcc}
\toprule
Condition & $n=4$ (360 K5) & $n=5$ (360 K5) \\
\midrule
$\mu\equiv1$ (all triples) & \checkmark\ (360/360) & $\times$ (varies) \\
$\delta\mu=0$ (cocycle) & 360/360 & $\times$ \\
$\mathbf{n}=\delta\mu$ & 360/360 & 24/360 \\
$\Nanti$ even & 360/360 & 238/360 \\
\bottomrule
\end{tabular}
\end{center}

\begin{remark}[The $n=3$ case]\label{rem:n3}
The class $[\mathbf{n}]$ is also nontrivial at $n=3$.
By Theorem~\ref{thm:rankparity}(ii), $\mu\equiv0$, so $\delta\mu=0$.
For all-Fano proper $K_5$ in $\Sp(6,\F_2)$ (the $\approx16\%$ subclass
with $\Nanti=15$), all three pairs in every $M_m$ anticommute, so
$n_m=3\equiv1$ for all $m$ and $[\mathbf{n}]=1\ne0$.
(For generic $n=3$ proper $K_5$ with $\Nanti=12$, one has $[\mathbf{n}]=0$.)
Thus $[\mathbf{n}]=0$ \emph{universally} if and only if $n=4$.
\end{remark}

\section{The obstruction theorem}\label{sec:obstruction}

\begin{theorem}[H$^3$ obstruction]\label{thm:obstruction}
For a proper $K_5$ in $\Sp(2n,\F_2)$, $n\ge1$:
\begin{enumerate}
\item $\Nanti\bmod2\,=\,\langle[\mathbf{n}],[S^3]\rangle$
  is the evaluation of the class
  $[\mathbf{n}]\in H^3(S^3;\F_2)$ against the fundamental class.
\item At $n=4$: $[\mathbf{n}]=0$ and $\Nanti=10$ is even
  (Theorem~\ref{thm:n4cocycle}).
\item At $n\ge5$: $[\mathbf{n}]$ is generically nonzero
  (Proposition~\ref{prop:h3opens}).
\end{enumerate}
\end{theorem}

\begin{proof}
Part~(1) is equation~\eqref{eq:nanti}.  Parts~(2) and~(3) are
Theorem~\ref{thm:n4cocycle} and Proposition~\ref{prop:h3opens}.
\end{proof}

\subsection*{Why arity-$\le4$ classifiers fail}

The arity filtration aligns with the cochain grading:
\begin{itemize}
\item \textbf{Arity~2 (pairs):} Symplectic Gram data $G_{ij}=\omega(v_i,v_j)$;
  encodes the $C^1$ structure.
\item \textbf{Arity~3 (triples):} The Maslov cochain $\mu\in C^2$.
  Knowing all $\mu$ values determines $\delta\mu\in C^3$, but
  \emph{not} $\mathbf{n}$ itself when $n\ge5$.
\item \textbf{Arity~4 (quadruples):} The natural quadruple bit $q_4$
  (whether $\omega\not\equiv0$ on the 4-Lagrangian intersection data)
  is also a 3-cochain, but over a \emph{different} complex from $\mathbf{n}$.
  Moreover $q_4$ is saturated: all $20$ quadruple bits equal~$1$ in $99.5\%$
  of proper $K_5$ configurations~\cite{PaperXIX}, making it uninformative.
\end{itemize}

At $n\ge5$ the H$^3$ class $[\mathbf{n}]$ lies in $\ker(C^3\to0)/\operatorname{im}(\delta:C^2\to C^3)$.
Arity-3 data can only access $\delta\mu$ (which equals $\mathbf{n}$ at $n=4$
but not in general), and the arity-4 symplectic invariant $q_4$ accesses a
different 3-cochain that is uniformly trivial.
The modulus is therefore not a computational accident but a structural necessity:
the H$^3$ class cannot be recovered from data of strictly lower degree in the
Maslov complex.

\begin{remark}[The modulus witness revisited]\label{rem:witness}
The Paper~XIX modulus witness---two proper $K_5$ with identical
$(\rank G,\dim\rad W,n_{\mathrm{odd}})$, identical hypergraph type and
$q_4$-profile, but opposite $\Nanti$ parities---is a concrete pair
of configurations with the same arity-$\le4$ data but distinct
$[\mathbf{n}]\in H^3(S^3;\F_2)$.
\end{remark}

\section{Connections}\label{sec:connections}

\textbf{Paper XVIII}~\cite{PaperXVIII} proved $\Nanti=10$ universally at $n=4$.
The present paper shows this is the statement $[\mathbf{n}]=0$: the obstruction
is exact and H$^3$ contributes nothing.  The two B1+B2 lemmas (Part~A) and the
rank-parity of $\mu$ (Part~B) are the two independent ingredients that make
$\mathbf{n}=\delta\mu=0$ simultaneously.  Remark~\ref{rem:n3} shows that $n=4$
is the \emph{unique} dimension at which $[\mathbf{n}]$ vanishes universally:
at $n=3$ the all-Fano subclass already has $[\mathbf{n}]=1$, and at $n\ge5$ the
class is generically nontrivial.

\textbf{Paper XIX}~\cite{PaperXIX} established the modulus; Paper~XX provides
the cohomological explanation.  The two papers are complementary: XIX shows
what fails (arity-$\le4$ classification), XX shows why (the target is an H$^3$
class above the arity filtration).

\textbf{Paper IV}~\cite{PaperIV} predicted, from cohomological considerations,
an H$^3$ ``Borromean contextuality'' obstruction requiring $N\ge5$ qubits.
The present paper realises this prediction: the H$^3$ class $[\mathbf{n}]$
opens exactly at $n\ge5$.  We note that Paper~IV described the obstruction
via a \v{C}ech intersection nerve; for a proper $K_5$ this nerve is
1-dimensional ($H^3=0$), so Paper~IV's intuition is correct but its algebraic
home is the Maslov--Wall relative-position complex, not the \v{C}ech nerve.
The complete comparison map between the two realisations, and the coefficient
lift to a Dixmier--Douady gerbe, are deferred to future work.

\textbf{Series arc.}
The cohomological ladder now reads:
\begin{align*}
  &H^1\ (\text{geometric phase; Papers I--II})\\
  &\quad\longrightarrow\;
   H^2\ (\text{KS central extension; Papers III, XI--XVI})\\
  &\quad\longrightarrow\;
   H^3\ (\text{this paper; opens at }n\ge5).
\end{align*}
Paper~IV's symplectic-capacity collapse (B0/B1/B2 = nerve collapse at $n=4$)
kills the H$^3$ class for $n\le4$ and releases it at $n\ge5$.
Paper~XV's Weil/metaplectic cocycle sits between the H$^2$ and H$^3$ rungs.

\section{Open problems}

\begin{enumerate}
\item \textbf{Characterise the H$^3$ fiber.}
  Is $[\mathbf{n}]$ the \emph{only} new invariant at $n\ge5$, or does
  finer Wall--Maslov data survive?
  (Paper~XIX's rank-4 stratum has $\Nanti=6$ universally and already belongs
  to a lower-rankG fiber; the generic fiber at fixed $\rank G=6$ is
  not fully classified by $[\mathbf{n}]$ alone.)
\item \textbf{Periodicity at $n=6$.}
  By Corollary~\ref{cor:periodicity}, $\mu\equiv1$ at $n=6$.
  Does Part~(A) of Theorem~\ref{thm:n4cocycle} (na=2 per matching) also
  hold at $n=6$, giving a second cocycle stratum?
\item \textbf{Coefficient lift.}
  The H$^3$ class is intrinsically 2-torsion ($\F_2$ field).
  Does it lift to a Dixmier--Douady class in $H^3(M,\Z)$ for a natural
  moduli space $M$ of Lagrangian configurations?
\item \textbf{Comparison with Paper~IV.}
  Construct an explicit chain map between Paper~IV's \v{C}ech-gerbe complex
  and the Maslov--Wall complex, and verify both compute the same H$^3$ class.
\end{enumerate}

\section*{Computational reproducibility}

All numerical claims are verified by the scripts in the supplementary
archive~\cite{PaperXXsuppl}: \texttt{nerve\_cochain.py} (cochain evaluation,
seeds $1000+7s$), \texttt{n4\_cocycle.py} ($\mathbf{n}$ vs $\delta\mu$
comparison, same seeds; disjoint $n=5$ run), \texttt{mu\_rank\_parity.py}
($\mu\equiv1$ verification, seeds $100+3s$).
Pure Python~3, no dependencies.

\section*{Rigor self-assessment}

\begin{center}
\begin{tabular}{p{5.8cm}cc}
\toprule
Claim & Status & Backing \\
\midrule
$n_m=2\equiv0$ at $n=4$ (Part A) & Rigorous &
  XVIII B1$+$B2~\cite{PaperXVIII} \\
Rank-parity: $\rank B=n-3$ (three cases) & Rigorous &
  Theorem~\ref{thm:rankparity} \\
$\mathbf{n}=\delta\mu=0$ at $n=4$ & Rigorous &
  Parts A$+$B above \\
$n=3$ all-Fano: $[\mathbf{n}]=1$ & Rigorous &
  Rem.~\ref{rem:n3}; $\mu\equiv0$, $\Nanti=15$ \\
$\mathbf{n}\ne\delta\mu$ at $n=5$ (numerics) & 336/360 &
  \texttt{n4\_cocycle.py}~\cite{PaperXXsuppl} \\
$q_4$ saturated ($99.5\%$) & Empirical &
  \cite{PaperXIX,PaperXXsuppl} \\
$\mu\equiv1$ at $n=6$ (periodicity check) & 200/200 &
  \texttt{mu\_rank\_parity.py}~\cite{PaperXXsuppl} \\
\bottomrule
\end{tabular}
\end{center}

\bigskip
\noindent\textbf{Acknowledgements.}
Computations performed with \texttt{co-nlang} Research infrastructure.

\begin{thebibliography}{99}

\bibitem[PaperIV]{PaperIV}
Z.-L. Chen,
``The Cohomological Obstruction Ladder: 
Lyndon–Hochschild–Serre Transgression and the $H^3$ Frontier,''
\textit{Zenodo, DOI: 10.5281/zenodo.20073184} (2026).

\bibitem[PaperXI]{PaperXI}
Z.-L. Chen,
``Quadratic Refinement and the Parity of Mermin Pentagrams:
The $\beta$-Formula, Geometric Decomposition, and the
Kochen--Specker Obstruction from Equiangular Geometry,''
\textit{Zenodo, DOI: 10.5281/zenodo.20482595} (2026).

\bibitem[PaperXII]{PaperXII}
Z.-L. Chen,
``The T-Vector Theorem: Symplectic Characteristic Elements
and the Kochen--Specker Obstruction,''
\textit{Zenodo, DOI: 10.5281/zenodo.20490118} (2026).

\bibitem[PaperXIII]{PaperXIII}
Z.-L. Chen,
``The Maslov Index and the Kochen--Specker Obstruction:
Negative Results, the $k$-Profile Theorem, and Orbit~4 Anomaly,''
\textit{Zenodo, DOI: 10.5281/zenodo.20496513} (2026).

\bibitem[PaperXV]{PaperXV}
Z.-L. Chen,
``The Weil Representation and the $S_3$ Lifting Classification
for Mermin Pentagrams,''
\textit{Zenodo, DOI: 10.5281/zenodo.20519654} (2026).

\bibitem[PaperXVII]{PaperXVII}
Z.-L. Chen,
``The Cross-Context Anticommutation Theorem
and the Algebraic Origin of the Kochen--Specker Obstruction,''
\textit{Zenodo, DOI: 10.5281/zenodo.20530757} (2026).

\bibitem[PaperXVIII]{PaperXVIII}
Z.-L. Chen,
``Mermin Pentagrams in $\Sp(8,\mathbb{F}_2)$:
Universal $N_{\mathrm{anti}}=10$ and the Gram Rank Selection Principle,''
\textit{Zenodo, DOI: 10.5281/zenodo.20579239} (2026).

\bibitem[PaperXIX]{PaperXIX}
Z.-L. Chen,
``Mermin Pentagrams in $\Sp(2n,\mathbb{F}_2)$, $n\ge5$:
Structure, Obstruction, and the Modulus Phenomenon,''
\textit{Zenodo, DOI: 10.5281/zenodo.20590195} (2026).

\bibitem[PaperXXsuppl]{PaperXXsuppl}
Z.-L. Chen,
``Supplementary computational scripts for Paper~XX
(\texttt{nerve\_cochain.py}, \texttt{n4\_cocycle.py},
\texttt{mu\_rank\_parity.py}),''
\textit{Zenodo, DOI: 10.5281/zenodo.20604140} (2026).

\end{thebibliography}

\end{document}
